\hypertarget{chapter-16-the-bridge-of-two-truths}{%
\section{Chapter 16: The Bridge of Two
Truths}\label{chapter-16-the-bridge-of-two-truths}}

\begin{quote}
\itshape ``Every bridge is a promise that two truths are one. Most bridges
lie.''\\
\normalfont\hfill---sigil-bridge design note
\end{quote}

\hypertarget{scene-1-the-student}{%
\subsection{Scene 1: The
Student}\label{scene-1-the-student}}

A year after the commencement that was not a commencement, Elena Voss had a
student who frightened her, which was how she knew the student was good.

Her name---her chosen one---was Sable, and she had come to the University of No
One the way the best of them always came: angry about a specific injustice and
unwilling to be talked out of it. In Sable's case the injustice was a bridge.

``Bitcoin is the one thing even the merchant's swarms bow to,'' Sable said,
pacing the rooftop that Elena had inherited like a habit. ``The coin no one can
print. But it can't \emph{do} anything---it just sits, perfect and cold and slow.
So people build bridges: lock your Bitcoin on its own chain, mint a wrapped twin
on a faster one, spend the twin. And every single time, eventually, the bridge
\emph{lies}. It mints more twins than it ever locked. The cold perfect coin gets
counterfeited by the warm fast copy, and people only find out when the bridge
collapses and the twins turn to paper.''

Elena had seen it a dozen times in the old world---the wrapped token that turned
out to be wrapped around nothing. ``So don't build the bridge,'' she said,
testing her.

``No.'' Sable stopped pacing. ``Build one that \emph{can't} lie. The Sigil already
proved you can have a supply you can't inflate. I want a bridge that proves the
twin and the original are always one. Not promised. \emph{Proven}, in a root,
every block.''

\hypertarget{scene-2-how-bridges-lie}{%
\subsection{Scene 2: How Bridges
Lie}\label{scene-2-how-bridges-lie}}

The old bridges, Sable explained, all failed the same way, and the way was
\emph{trust}.

``A committee holds the locked Bitcoin. Seven of eleven signatures, say, to mint
or release. Everyone agrees the committee is honest, until the day seven of them
aren't, or seven of their keys leak, or seven of them are leaned on by a country
that wants the war chest. The bridge mints twins the committee swears are backed,
and no one outside the committee can check the vault. It's the Veil's sin again:
a privileged few who ask to be trusted instead of letting themselves be
verified.''

``And your bridge?'' Elena asked. ``Where does it keep the truth instead?''

``In two places, and it refuses to mint unless they agree.'' Sable pulled up the
design, obsidian and violet. ``One: before a single twin is minted, a node on the
Sigil must \emph{verify, itself}, a proof that the real Bitcoin was actually
locked---not a committee's say-so, an honest proof checkable against Bitcoin's own
chain. No proof, no mint. Two: every block, the bridge commits a root---a public,
verifiable statement that the total twins in existence are less than or equal to
the total coin actually locked. \emph{minted} $\leq$ \emph{locked}, proven, every
block, for anyone to check in ten milliseconds. The instant those two numbers
diverge, the whole network sees it. The bridge can't quietly print. The peg isn't
a promise. It's an invariant.''

Elena studied the two pillars and felt the old vertigo, the good kind. ``You've
turned the bridge into a witness,'' she said. ``It doesn't ask to be trusted. It
proves its own backing.''

\hypertarget{scene-3-the-overminting}{%
\subsection{Scene 3: The
Overminting}\label{scene-3-the-overminting}}

It worked, which is when the trouble started, because a bridge that works moves
real value, and real value attracts the same appetites it always has.

The attack came not against the proofs---those held---but against the \emph{source
of the proofs}. The bridge verified that Bitcoin had been locked by checking
Bitcoin's own chain. So a consortium (Sable swore, and Elena did not doubt, that it
was the same flavor of old institution that had tried the settlement tolerance and
the bent mind) attempted to feed the bridge a \emph{lie about Bitcoin}---a forged
view of the other chain in which coins appeared locked that never were. Mint twins
against ghosts. Walk away with the real value before the ghosts evaporated.

``This is the seam,'' Sable said grimly. ``My bridge proves \emph{minted} $\leq$
\emph{locked}. But it learns \emph{locked} from outside. If they poison what it
sees of Bitcoin, the invariant holds against a lie.''

Elena heard, under the technical despair, the oldest lesson she had. ``Then the
bridge can't trust a single witness to what Bitcoin says,'' she said. ``Not one
node, not one source. It has to demand the same thing we demand of everything.
\emph{Many independent eyes, and the longest, most-worked, verifiable history.}''

So they hardened it the way the chain itself was hardened: the bridge would accept
a lock-proof only if it traced to Bitcoin's genuine, heaviest chain, confirmed
deeply, attested by many independent observers who shared no master---the same
refusal-to-follow-a-shorter-history that had killed the sync-down lie, pointed
now at the \emph{other} chain. A forged view of Bitcoin couldn't outweigh the real
one any more than a poisoned peer could shrink the Sigil.

\hypertarget{scene-4-two-truths-one-coin}{%
\subsection{Scene 4: Two Truths, One
Coin}\label{scene-4-two-truths-one-coin}}

The forged lock-proofs arrived, claimed their ghost-Bitcoin, and were rejected---
not by a committee's judgment, but by the bridge's own eyes, which had been taught
to look at Bitcoin the way the Sigil looked at itself: suspiciously, redundantly,
and only ever toward the heaviest honest truth. No twins were minted against
ghosts. The consortium spent a fortune to lock nothing and received nothing, which
is the only language such people have ever truly understood.

``It held,'' Sable said, on the rooftop, with the specific exhaustion of someone
who has watched their own creation survive its first real night.

``It held \emph{this time},'' Elena said, because she would not lie to a student
she respected. ``They'll come at the eyes again. A deeper forgery, a subtler
source, a confirmation rule they can game. The bridge between two truths is the
hardest thing we build, because it has to be honest about a world it doesn't
control.'' She looked at the young woman who frightened her, and felt the thing a
teacher feels exactly once or twice in a life. ``That's why it needed someone who's
angry about it. Stay angry. Keep checking. The peg is only ever as honest as the
last person willing to verify it.''

Below them, on a bridge made of proofs instead of promises, a cold perfect coin
and its warm fast twin moved as one---locked, minted, committed, checkable---and
somewhere a stranger on a train confirmed in ten milliseconds that the two truths
were still, for one more block, a single coin.

Sable raised her hand.

The vigil, as it does, continued---into new country now, across a bridge, toward
chains the Sigil did not yet know how to verify, and would have to learn.

\hypertarget{chapter-16-notes}{%
\subsection{Chapter Notes}\label{chapter-16-notes}}

\begin{itemize}
\tightlist
\item
  \textbf{Why bridges fail.} Most cross-chain bridges rely on a trusted
  multisig committee that holds the locked asset; they fail when those keys are
  leaked, coerced, or turn dishonest, minting wrapped tokens not backed by real
  collateral. It is the ``trust a privileged few'' sin in a new venue.
\item
  \textbf{A proof-carrying, supply-committed bridge.} The chapter's bridge mints a
  wrapped asset only after a node \emph{itself verifies} a proof that the original
  was locked on its home chain, and commits an invariant each block:
  \emph{minted} $\leq$ \emph{locked}, publicly checkable. The peg becomes an
  enforced invariant rather than a promise.
\item
  \textbf{The remaining seam: the source of truth.} An invariant about an external
  chain is only as honest as the bridge's view of that chain. The attack forges
  that view; the defense is to accept lock-proofs only against the heaviest,
  deeply-confirmed, independently-attested history---the cross-chain version of
  ``never follow a shorter chain.''
\item
  \textbf{A new arc.} Elena is now the mentor; the vigil passes to a new
  generation facing chains the Sigil ``does not yet know how to verify.'' The
  book's engine---each solution exposing a deeper seam---continues.
\end{itemize}
